\exercisesection{DIMENSION OF GRADED RINGS}

\begin{enumerate}
\def\labelenumi{\arabic{enumi})}
\tightlist
\item
  Let A be a ring, and let p be a graded prime ideal of A{[}T{]}, and \(p_{0}\) the prime ideal \(p \cap A\) of A.
\end{enumerate}

\begin{enumerate}
\def\labelenumi{\alph{enumi})}
\item
  If \(T \in p\) , then \(p = p_{0} + T \cdot A[T]\) .
\item
  If T \(\notin\) p, then \(\mathfrak{p} = \mathfrak{p}_0\) . A{[}T{]}.
\item
  Deduce that, in the case \(a)\) we have \(\operatorname{htgr}(\mathfrak{p}) = \operatorname{ht}(\mathfrak{p}_0) + 1\) , and in the case \(b)\) , one has \(\operatorname{htgr}(\mathfrak{p}) = \operatorname{ht}(\mathfrak{p}_0)\) .
\item
  Show that \(\dim \operatorname{gr}(\mathbf{A}[\mathbf{T}]) = \dim (\mathbf{A}) + 1\) .
\item
  Deduce from this an example of a graded ring B such that dim(B) ≠ dimgr(B) (p.84, exerc. 7).
\end{enumerate}

\begin{enumerate}
\def\labelenumi{\arabic{enumi})}
\setcounter{enumi}{1}
\tightlist
\item
  Let A be a ring, and \(\mathcal{F} = (\mathcal{F}_i)_{i \in \mathbf{Z}}\) a decreasing filtration of A such that \(\mathcal{F}_i = \mathbf{A}\) for \(i \leqslant 0\). Let \(\mathcal{A}\) be the subring \(\bigoplus_{i \in \mathbf{Z}} \mathcal{F}_i \cdot v^{-i}\) of \(\mathbf{A}[v, v^{-1}]\), and let \(\mathrm{gr}_{\mathcal{F}}(\mathbf{A})\) be the associated graded ring
\end{enumerate}

of the filtered ring (A, F).

\begin{enumerate}
\def\labelenumi{\alph{enumi})}
\item
  Show that the homomorphism \(f: \mathcal{A} \to \mathrm{gr}_{\mathcal{F}}(\mathbf{A})\) defined by \(f\left(\sum_{i} a_i v^{-i}\right) = \sum_{i} \overline{a}_i\), where \(\overline{a}_i\) is the class of \(a_i\) in \(\mathcal{F}_i / \mathcal{F}_{i+1}\), induces an isomorphism of \(\mathcal{A}/v. \mathcal{A}\) with \(\mathrm{gr}_{\mathcal{F}}(\mathbf{A})\).
\item
  Show that for every invertible element \(v_{0}\) of A, the homomorphism \(e_{v_{0}}: \mathcal{A} \to \mathbf{A}\) defined by
\end{enumerate}

\[
e _ {v _ {0}} (\sum a _ {i} v ^ {- i}) = \sum a _ {i} v _ {0} ^ {- i}
\]

induces an isomorphism of \(\mathcal{A}/(v - v_{0})\mathcal{A}\) with A.

\begin{enumerate}
\def\labelenumi{\alph{enumi})}
\setcounter{enumi}{2}
\item
  Show that \(\mathcal{A}/\mathrm{A}[v]\) is a torsion \(\mathrm{A}[v]\)-module.
\item
  Now assuming that A contains a field k, show that the composite morphism
\end{enumerate}

\[
k [ v ] \rightarrow \mathrm{A} [ v ] \rightarrow \mathcal {A} \quad (\text {where } \mathrm{A} [ v ] = \bigoplus_ {i \leqslant 0} \mathcal {F} _ {i}. v ^ {- i})
\]

makes \(\mathcal{A}\) into a k{[}v{]}-torsion-free module. Deduce that \(\mathcal{A}\) is a flat k{[}v{]}-module.

\begin{enumerate}
\def\labelenumi{\alph{enumi})}
\setcounter{enumi}{4}
\tightlist
\item
  Suppose that A is local, with maximal ideal m, and contains a field k such that A is the direct sum of k and m (as an additive group). Show that there exists an ideal S of A such that \(\mathcal{A}\) is the direct sum of k{[}v{]} and S (we have set \(\mathcal{F}_{i} = m^{i}\) for every \(i \in \mathbf{Z}\)).
\end{enumerate}

\begin{enumerate}
\def\labelenumi{\arabic{enumi})}
\setcounter{enumi}{2}
\tightlist
\item
  Let us resume the notation of the previous exercise. Let M be an A-module, equipped with a filtration \((\mathbf{K}_i)_{i\in\mathbf{Z}}\) such that \(\mathbf{K}_i = \mathbf{M}\) for \(i \leqslant 0\).
\end{enumerate}

\begin{enumerate}
\def\labelenumi{\alph{enumi})}
\tightlist
\item
  Show that the graded A-module \(\mathcal{M} = \sum_{i\in \mathbf{Z}}\mathrm{K}_{i}v^{-i}\) has a natural structure of \(\mathcal{A}\) -graded module.
\end{enumerate}



\begin{enumerate}
\def\labelenumi{\alph{enumi})}
\setcounter{enumi}{1}
\item
  Show that the graded A-algebra \(\mathcal{A}\) is generated by its elements of degree 1 and -1 if and only if there exists an ideal q of A such that \(\mathcal{F}_i = \mathfrak{q}^i\) ( \(i \in \mathbf{Z}\) ).
\item
  Assuming this last condition is satisfied, show that the filtration \((\mathbf{K}_i)_{i\in \mathbb{Z}}\) is q-good if and only if the \(\mathcal{A}\) -graded module \(\mathcal{M}\) is of finite type.
\item
  Show that if the filtration \((\mathbf{K}_{i})_{i\in\mathbf{Z}}\) is q-good, the A-module M is v-torsion-free, and conversely, given a finitely generated A-module M without v-torsion, one can equip it with a q-good filtration such that the associated A-module is M.
\end{enumerate}

\begin{enumerate}
\def\labelenumi{\arabic{enumi})}
\setcounter{enumi}{3}
\tightlist
\item
  Let us take up the notations and hypotheses of exercise 2. We say that the filtration \(\mathcal{F}\) is a filtration by quasi-powers if there exists an integer \(k > 0\) such that
\end{enumerate}

\[
\mathcal {F} _ {n} = \sum_ {i = 1} ^ {k} \mathcal {F} _ {i}. \mathcal {F} _ {n - i} \quad \text { for } \quad n \geqslant 1.
\]

Show that if the ring A is Noetherian, the following conditions are equivalent:

\begin{enumerate}
\def\labelenumi{(\roman{enumi})}
\item
  The filtration \(\mathcal{F}\) is a filtration by quasi-powers.
\item
  The ring \(\mathcal{A}\) is Noetherian.
\item
  The ideal \(\mathcal{N} = \sum_{i > 1} v^{-i}\mathcal{F}_i\) of \(\mathcal{A}\) is of finite type.
\item
  There exists an integer \(k \geqslant 1\) such that \(\mathcal{F}\) is the smallest filtration of A for the order relation “\(\mathcal{H}_n \subset \mathfrak{G}_n\) for every \(n\)”, whose first terms are \(\mathcal{F}_1, ..., \mathcal{F}_k\).
\item
  The A-algebra \(\mathcal{A}\) is of finite type.
\end{enumerate}

\begin{enumerate}
\def\labelenumi{\arabic{enumi})}
\setcounter{enumi}{4}
\tightlist
\item
  Let \(\rho: A \to B\) be an injective homomorphism of Noetherian rings making \(B\) an \(A\)-algebra of finite type. We call the transcendence degree of \(B\) over \(A\), and denote it by \(d_A(B)\), the largest integer \(d\) such that there exists an injective homomorphism of \(A\)-algebras \(b: A[X_1, \ldots, X_d] \to B\). a) Show that if \(\sigma: B \to C\) satisfies the same hypotheses as \(\rho\), one has the inequality
\end{enumerate}

\[
\mathrm{d} _ {\mathbf {A}} (\mathbf {C}) \geqslant \mathrm{d} _ {\mathbf {A}} (\mathbf {B}) + \mathrm{d} _ {\mathbf {B}} (\mathbf {C}).
\]

\begin{enumerate}
\def\labelenumi{\alph{enumi})}
\setcounter{enumi}{1}
\item
  Show that if A and B are integral domains, \(d_{\mathbf{A}}(\mathbf{B})\) is equal to the transcendence degree of the field of fractions of B over that of A. Moreover, if in a) one assumes C is an integral domain, one has the equality \(d_{\mathbf{A}}(\mathbf{C}) = d_{\mathbf{A}}(\mathbf{B}) + d_{\mathbf{B}}(\mathbf{C})\).
\item
  Show that if B is a subalgebra of A{[}v, \(v^{-1}\) {]} containing A{[}v{]}, where v is an indeterminate, we have \(d_{\mathbf{A}}(\mathbf{B}) = 1\) .
\end{enumerate}

6) Let \(\rho: A \to B\) be an injective homomorphism of Noetherian rings, making \(B\) an \(A\)-algebra of finite type. Let \(q\) be a prime ideal of \(B\) distinct from \(B\), and set \(p = \rho^{-1}(q)\). a) Show that, with the notations of the preceding exercise, one has the inequality

\[
(*) \quad \mathrm{ht} (\mathfrak {q}) + \mathrm{d} _ {\mathrm{A} / \mathfrak {p}} (\mathrm{B} / \mathfrak {q}) \leqslant \mathrm{ht} (\mathfrak {p}) + \mathrm{d} _ {\mathrm{A}} (\mathrm{B}).
\]

(One may proceed by induction on \(d_{A}(B)\) .)

\begin{enumerate}
\def\labelenumi{\alph{enumi})}
\setcounter{enumi}{1}
\item
  Show that if the A-algebra B is generated by d elements and \(\mathrm{d}_{\mathbf{A}}(\mathbf{B}) = d\) one has equality in (*).
\item
  Show that if A is an integral domain and the polynomial ring A{[}T \(_{1}\) , \ldots, T \(_{n}\) {]} is catenary for every n, equality holds in (*) for every integral A-algebra B of finite type and every prime ideal q of B distinct from B.
\end{enumerate}

\begin{enumerate}
\def\labelenumi{\arabic{enumi})}
\setcounter{enumi}{6}
\tightlist
\item
  Let A be a Noetherian ring and \(\mathcal{F} = (\mathcal{F}_i)_{i\in \mathbf{Z}}\) a decreasing filtration of A such that \(\mathcal{F}_0 = \mathbf{A}\), \(\mathcal{F}_1 \neq \mathbf{A}\), and the A-algebra \(\mathbf{B} = \sum_{i=1}^{\infty} \mathcal{F}_i v^{-i}\) is of finite type.
\end{enumerate}

We propose to show that we have \(\dim(B) = \dim(A) + 1\) and to use this fact to generalize the result of § 6, no. 3, corollary of prop. 5. One will use the results of the two preceding exercises.

\begin{enumerate}
\def\labelenumi{\alph{enumi})}
\tightlist
\item
  Let M be a maximal ideal of B. Show that we have
\end{enumerate}

\[
\operatorname{ht} (\mathfrak {M}) \leqslant \dim (\mathrm{A}) + 1,
\]

and deduce from it the inequality

\[
\dim (\mathbf {B}) \leqslant \dim (\mathbf {A}) + 1.
\]

\begin{enumerate}
\def\labelenumi{\alph{enumi})}
\setcounter{enumi}{1}
\tightlist
\item
  Using the fact that every element of the A{[}v{]}-module A{[}v, \(v^{-1}\) {]}/A{[}v{]} is annihilated by a power of v, show that we have
\end{enumerate}

\[
\dim (\mathrm{A} [ v, v ^ {- 1} ]) = \dim (\mathrm{A} [ v ]) = \dim (\mathrm{A}) + 1,
\]

and deduce that we have \(\dim(\mathbf{B}) \geqslant \dim(\mathbf{A}) + 1\) , whence finally \(\dim(\mathbf{B}) = \dim(\mathbf{A}) + 1\) (one may use the fact that the ring \(\mathbf{A}[v, v^{-1}]\) is a ring of fractions of \(\mathbf{B}\) ). c) Show that the minimal prime ideals of \(\mathbf{B}\) are the ideals of the form:

\[
\sum_ {i \in \mathbf {Z}} \left(\mathfrak {p} \cap \mathcal {F} _ {i}\right) v ^ {- i}
\]

where p is a minimal prime ideal of A.

\begin{enumerate}
\def\labelenumi{\alph{enumi})}
\setcounter{enumi}{3}
\item
  Show that one has \(v\mathbf{B} \neq \mathbf{B}\) and that \(v\) does not belong to any of the minimal prime ideals of B. Deduce the inequality \(\dim(\mathbf{B}/v\mathbf{B}) \leqslant \dim(\mathbf{A})\) (one may use § 1 and § 3, no. 3, cor. 1 of prop. 4).
\item
  Suppose now that there exists a maximal ideal m of A containing \(\mathcal{F}_1\) and such that ht(m) = dim(A). Show that in this case we have dim(B/vB) ≥ dim(A) and therefore dim(B/vB) = dim(A). (One may reduce to the case where A is a local ring with maximal ideal m, show that the ideal M of B generated by v and \(\sum_{i\in\mathbf{Z}}(\mathcal{F}_i\cap m)v^{-i}\) is a maximal ideal of B of height dim(A) + 1, and localize B at M.)
\end{enumerate}





\begin{enumerate}
\def\labelenumi{\alph{enumi})}
\setcounter{enumi}{5}
\tightlist
\item
  If A is a Noetherian local ring, and F is a decreasing filtration of A such that \(F_{0} = A\) , \(F_{1} \neq A\) , we have the equality
\end{enumerate}

\[
\dim (\operatorname{gr} (A)) = \dim (A).
\]

\begin{enumerate}
\def\labelenumi{\arabic{enumi})}
\setcounter{enumi}{7}
\tightlist
\item
  Let A be a Noetherian ring, and let B be a graded sub-A-algebra of the polynomial ring A{[}T{]}.
\end{enumerate}

\begin{enumerate}
\def\labelenumi{\alph{enumi})}
\tightlist
\item
  Show that we have the inequalities
\end{enumerate}

\[
\dim (\mathbf {A}) \leqslant \dim (\mathbf {B}) \leqslant \dim (\mathbf {A}) + 1.
\]

\begin{enumerate}
\def\labelenumi{\alph{enumi})}
\setcounter{enumi}{1}
\tightlist
\item
  Let \(B_{i}\) be the set of \(b \in A\) such that \(bT^{i} \in B\). Show that if we have \(\dim(B) = \dim(A)\), there exists a positive integer k such that for all \(i \geqslant k\), \(B_{i}\) is contained in the intersection of the minimal prime ideals of A.
\end{enumerate}

\begin{enumerate}
\def\labelenumi{\arabic{enumi})}
\setcounter{enumi}{8}
\item
  Let A be a ring and M a finitely presented A-module. For every \(\mathfrak{p} \in \operatorname{Spec}(A)\), we set \(\kappa(\mathfrak{p}) = A_{\mathfrak{p}} / \mathfrak{p} \cdot A_{\mathfrak{p}}\). Show that for all \(k \in \mathbf{N}\), the set of \(\mathfrak{p} \in \operatorname{Spec}(A)\) such that \(\dim_{\kappa(\mathfrak{p})} (M \otimes_A \kappa(\mathfrak{p})) \geqslant k\) is a closed subset of \(\operatorname{Spec}(A)\). (One may take a presentation \(\mathrm{A}^m \xrightarrow{u} \mathrm{A}^n \longrightarrow \mathrm{M} \longrightarrow 0\) of M and consider the minors of the matrix of \(u\).)
\item
  Let \(\rho: A \to B\) be a ring homomorphism. For every \(\mathfrak{p} \in \operatorname{Spec}(B)\), we set
\end{enumerate}

\[
\mathrm{d} (\mathfrak {p}) = \dim_ {\mathfrak {p}} \left(\mathbf {B} \otimes_ {\mathbf {A}} \kappa \left(\rho^ {- 1} (\mathfrak {p})\right)\right).
\]

The aim of this exercise is to prove that when B is a finitely generated A-algebra, the map \(p \mapsto d(p)\) from \(\operatorname{Spec}(B)\) to \(\mathbf{R}\) is upper semicontinuous (“Chevalley’s theorem”).

Denote by \(\delta \subset \mathbf{B} \otimes_{\mathbf{A}} \mathbf{B}\) the kernel of the canonical A-homomorphism \(\mathbf{B} \otimes_{\mathbf{A}} \mathbf{B} \to \mathbf{B}\). For every integer \(r\), set

\[
\mathrm{P} _ {\mathbf {B} / \mathbf {A}} ^ {r} = (\mathbf {B} \otimes_ {\mathbf {A}} \mathbf {B}) / \delta^ {r}.
\]

The B ⊗\_A B-algebra P \(_{B/A}^{r}\) is called the algebra of principal parts of order r. We consider B ⊗\_A B as a B-algebra via the homomorphism b ↦ b ⊗ 1, so that P \(_{B/A}^{r}\) is, for every r, a B-algebra and the canonical quotient homomorphisms P \(_{B/A}^{r+1}\) → P \(_{B/A}^{r}\) define a projective system of B-algebras.

\(^{1}\) This is the dimension as a vector space over the field \(\kappa(\mathfrak{p})\).

\begin{enumerate}
\def\labelenumi{\alph{enumi})}
\tightlist
\item
  Let \(p: \mathbf{A} \to \mathbf{A}'\) be a ring homomorphism. Set \(\mathbf{B}' = \mathbf{B} \otimes_{\mathbf{A}} \mathbf{A}'\) and denote by \(q: \mathbf{B} \to \mathbf{B}'\) the canonical homomorphism. Let \(r\) be an integer. Show that the unique homomorphism of \(\mathbf{B}'\)-algebras
\end{enumerate}

\[
\mathbf {P} _ {\mathbf {B} / \mathbf {A}} ^ {r} \otimes_ {\mathbf {B}} \mathbf {B} ^ {\prime} \rightarrow \mathbf {P} _ {\mathbf {B} ^ {\prime} / \mathbf {A} ^ {\prime}} ^ {r}
\]

which associates to the class of \((b_{1} \otimes b_{2}) \otimes 1\) (\(b_{i} \in \mathbf{B}\)) the class of \(q(b_{1}) \otimes q(b_{2})\) is an isomorphism. b) Let S be a multiplicative subset of B. Show that the canonical homomorphism

\[
\mathrm{S} ^ {- 1} \mathrm{P} _ {\mathrm{B/A}} ^ {r} \rightarrow \mathrm{P} _ {\mathrm{S} ^ {- 1} \mathrm{B/A}} ^ {r}
\]

is an isomorphism.

\begin{enumerate}
\def\labelenumi{\alph{enumi})}
\setcounter{enumi}{2}
\item
  Assume that B is a finitely presented A-algebra, that is, isomorphic to a quotient of \(\mathrm{A}[\mathrm{T}_{1},\ldots,\mathrm{T}_{n}]\) by a finitely generated ideal. Show that for every \(r\in\mathbf{N}\) , \(\mathbf{P}_{\mathrm{B}/\mathrm{A}}^{r}\) is a finitely presented B-module. For every \(\mathfrak{p}\in\operatorname{Spec}(\mathbf{B})\) , we set \(\mathrm{P}_{\mathrm{B}/\mathrm{A}}^{\infty}[\mathfrak{p}]=\lim_{r}\mathrm{P}_{\mathrm{B}/\mathrm{A}}^{r}\otimes_{\mathrm{B}}\kappa(\mathfrak{p})\) . Show that \(\mathrm{P}_{\mathrm{B}/\mathrm{A}}^{\infty}[\mathfrak{p}]\) is a \(\kappa(\mathfrak{p})\) -algebra that is local and complete with residue field \(\kappa(\mathfrak{p})\) . Let \(k\in\mathbf{N}\) . Show that the set of \(\mathfrak{p}\in\operatorname{Spec}(\mathbf{B})\) such that \(\dim(\mathrm{P}_{\mathrm{B}/\mathrm{A}}^{\infty}[\mathfrak{p}])\geqslant k\) is a closed subset of \(\operatorname{Spec}(\mathbf{B})\) . (For every \(r\in\mathbf{N}\) , let \(\mathrm{F}_{r,k}\) the set of \(\mathfrak{p}\) such that \(\dim_{\kappa(\mathfrak{p})}(\mathrm{P}_{\mathrm{B}/\mathrm{A}}^{r}\otimes\kappa(\mathfrak{p}))\geqslant\binom{r+k-1}{r-1}\) . We shall show that \(\mathrm{F}_{r,k}\) is closed (Exercise 9) and consequently that \(\Phi_{k}=\bigcap_{r}\mathrm{F}_{r,k}\) is closed. One will then apply § 6, no. 3, Cor. 2 to Prop. 6.)
\item
  Let m be an ideal of B and \(\mu\subset(\mathrm{B/m})\otimes_{\mathrm{A}}\mathrm{B}\) the kernel of the canonical mapping \((\mathrm{B/m})\otimes_{\mathrm{A}}\mathrm{B}\to\mathrm{B/m}\). Show that we have a canonical isomorphism
\end{enumerate}

\[
\mathrm{P} _ {\mathrm{B/A}} ^ {r} \otimes_ {\mathrm{B}} (\mathrm{B/m}) \rightarrow ((\mathrm{B/m}) \otimes_ {\mathrm{A}} \mathrm{B}) / \mu^ {r}.
\]

Assume that A is a field, that B is a finitely generated A-algebra, and that m is a maximal ideal of B. Show that \(P_{B/A}^{\infty}[m]\) is isomorphic to the completed localization of \((B/m)\otimes_{A}B\) at the maximal ideal \(\mu\). Deduce that one has \(\dim(P_{B/A}^{\infty}[m])=\dim_{m}(B)\). (Apply the zero theorem and th. 1 of § 2, no. 3. Note that the extension B/m of A is contained in a finite Galois extension of a finite purely inseparable extension of A.)

\begin{enumerate}
\def\labelenumi{\alph{enumi})}
\setcounter{enumi}{4}
\tightlist
\item
  Assume that B is a finitely generated A-algebra. Let \(p \in \operatorname{Spec}(B)\). Show that we have
\end{enumerate}

\[
\dim (\mathrm{P} _ {\mathrm{B/A}} ^ {\infty} [ \mathfrak {p} ]) = \dim \left(\mathrm{B} \otimes_ {\mathrm{A}} \kappa (\rho^ {- 1} (\mathfrak {p}))\right).
\]

First note that for every \(r \in \mathbf{N}\), one has \(\mathbf{P}_{\mathbf{B}/\mathbf{A}}^{\mathbf{r}} \otimes_{\mathbf{B}} \kappa(\mathfrak{p}) = (\mathbf{P}_{\mathbf{B}/\mathbf{A}}^{\mathbf{r}} \otimes_{\mathbf{B}} \mathbf{B}') \otimes_{\mathbf{B}'} \kappa(\mathfrak{p})\) by setting \(\mathbf{B}' = \mathbf{B} \otimes_{\mathbf{A}} \kappa(\rho^{-1}(\mathfrak{p}))\), and hence, by virtue of \(a\) ), we have \(\mathbf{P}_{\mathbf{B}/\mathbf{A}}^{\infty}[\mathfrak{p}] = \mathbf{P}_{\mathbf{B}' / \kappa(\rho^{-1}(\mathfrak{p}))}^{\infty}[\mathfrak{p}\mathbf{B}']\). One thus reduces to proving, in the case where A is a field, the equality

\[
\dim \left(P _ {B / A} ^ {\infty} [ p ]\right) = \dim_ {p} (B).
\]

Observe then that in the case where p is maximal, this equality follows from d). In the general case, let \(q_{1}, \ldots, q_{l}\) be the minimal prime ideals of B contained in p. Show, using the Nullstellensatz, that there exists a dense open set U of \(V(p)\) such that for every maximal ideal m belonging to U, the minimal prime ideals contained in m are contained in U. Using c), show that one may, by shrinking U, assume moreover that for every \(p' \in U\), one has \(\dim(\mathrm{P}_{\mathrm{B/A}}^{\infty}[\mathfrak{p}]) = \dim(\mathrm{P}_{\mathrm{B/A}}^{\infty}[\mathfrak{p}'])\). Conclude.)

\begin{enumerate}
\def\labelenumi{\alph{enumi})}
\setcounter{enumi}{5}
\tightlist
\item
  Assume that B is a finitely generated A-algebra. Show that the function \(\mathfrak{p} \mapsto \mathrm{d}(\mathfrak{p}) = \dim_{\mathfrak{p}} (\mathrm{B} \otimes_{\mathbf{A}} \kappa(\rho^{-1}(\mathfrak{p})))\) is upper semicontinuous. (When B is of finite presentation, use \(e\) and \(c\). In the general case, B is of the form \(\mathrm{A}[\mathrm{T}_1, ..., \mathrm{T}_n]/\mathfrak{J}\). Let \(\mathfrak{J}_{\alpha}\) be the family of finitely generated ideals of \(\mathrm{A}[\mathrm{T}_1, ..., \mathrm{T}_n]\) contained in \(\mathfrak{J}\), ordered by inclusion, and set \(\mathrm{B}_{\alpha} = \mathrm{A}[\mathrm{T}_1, ..., \mathrm{T}_n]/\mathfrak{J}_{\alpha}\). For every \(\alpha\), one has \(\operatorname{Spec}(\mathbf{B}) \subset \operatorname{Spec}(\mathrm{B}_{\alpha}) \subset \operatorname{Spec}(\mathrm{A}[\mathrm{T}_1, ..., \mathrm{T}_n])\), and \(\operatorname{Spec}(\mathbf{B}) = \bigcap \operatorname{Spec}(\mathrm{B}_{\alpha})\). For every \(\mathfrak{p} \in \operatorname{Spec}(\mathbf{B})\), set
\end{enumerate}

\[
\mathrm{d} _ {\alpha} (\mathfrak {p}) = \dim_ {\mathfrak {p}} \left(\mathrm{B} _ {\alpha} \otimes \kappa \left(\rho^ {- 1} (\mathfrak {p})\right)\right).
\]

Show that \(d_{\alpha}(p) \geqslant d(p)\) and that \(d(p) = \inf_{\alpha} d_{\alpha}(p)\) by virtue of the Noetherian character of \(\text{Spec}(\kappa(\rho^{-1}(p))[T_1, \ldots, T_n])\). Deduce that \(p \mapsto d(p)\) is the infimum of a filtered decreasing family of upper semicontinuous functions.)

\begin{enumerate}
\def\labelenumi{\arabic{enumi})}
\setcounter{enumi}{10}
\tightlist
\item
  Let \(\rho: \mathbf{A} \to \mathbf{B}\) be a ring homomorphism, and let \(^a\rho: \operatorname{Spec}(\mathbf{B}) \to \operatorname{Spec}(\mathbf{A})\) be the corresponding map.
\end{enumerate}

\begin{enumerate}
\def\labelenumi{\alph{enumi})}
\item
  Show that for every \(\mathfrak{q} \in \operatorname{Spec}(\mathbf{A})\), \(\operatorname{Spec}(\mathbf{B} \otimes_{\mathbf{A}} \kappa(\mathfrak{q}))\) is identified with \(^a\rho^{-1}(\mathfrak{q})\).
\item
  Suppose that B is a finitely generated A-algebra. Show that p is an isolated point in its fiber \(^a\rho^{-1}(^a\rho(p))\) if and only if \(\dim_p(B \otimes_A \kappa(^a\rho(p)))\) is zero.
\item
  From exercise 10, deduce that under the hypotheses of \(b\), the set of points of \(\operatorname{Spec}(\mathbf{B})\) which are isolated in their fiber is an open subset of \(\operatorname{Spec}(\mathbf{B})\).
\end{enumerate}

\begin{enumerate}
\def\labelenumi{\arabic{enumi})}
\setcounter{enumi}{11}
\tightlist
\item
  Let H be a finitely generated graded algebra, such that \(H_{0}\) is a field, and \((\xi_{a})_{a\in\mathbf{A}}\) homogeneous elements of H that generate H. Suppose there exists a family of strictly positive integers \((d_{i})_{i\in\mathbf{I}}\) such that \(P_{H}=\prod(1-T^{d_{i}})^{-1}\). Denote by \(\delta_{a}>0\) the degree of \(\xi_{a}\).
\end{enumerate}

\begin{enumerate}
\def\labelenumi{\alph{enumi})}
\item
  Show that there exists an injection \(\varphi: I \to A\) such that for every \(i\), \(d_i\) divides \(\delta_{\varphi(i)}\). (One may use § 4, no. 2, th. 1.)
\item
  Show that one has \(\inf_{a \in A} \delta_a \leqslant \inf_{i \in I} d_i\).
\item
  From this, deduce that if the \(\delta_{a}\) are equal to one another, the graded algebra H is regular.
\item
  We assume that A has two elements \(a\) and \(a'\) such that \(\delta_{a} < \delta_{a'}\) and that \(\delta_{a'}\) is not a multiple of \(\delta_{a}\). Show that the graded algebra H is regular.
\end{enumerate}

